\scanpage{45}
% archive sheet number: 44
\mspage{1}
\begingroup\csname tagsleft@true\endcsname
1. Opération de groupes \(\simeq\mathbb D_2\) sur la droite projective, et configuration octaédrale.

Soit \(V\) un groupe qui est un espace vectoriel de rang \(2\) sur \(\mathbb F_2\), \deleted{\unsure{?}} et considérons un\deleted{e droite pro} fibré en droites projectives \(X\) sur un schéma de base \(S\). On s'intéresse aux opérations, fidèles sur chaque fibre, de \(V\) sur \(X\), i.e. aux monomorphismes
\[
V_S\lhook\joinrel\longrightarrow\underline{\operatorname{Aut}}_S(X)=G \tag*{(1)}
\]
Choisissons une base \(e_1,e_2\) de \(V\) sur \(\mathbb F_{\unsure{1|2}}\), alors la donnée d'un hom
\[
V_S\longrightarrow\underline{\operatorname{Aut}}_S(X)\qquad\text{i.e.}\quad V\longrightarrow\operatorname{Aut}_S(X)
\]
qui soit mono sur \(\mathbb F_2 e_1\) et sur \(\mathbb F_2 e_2\), équivaut à celle de \inserted{\unsure{autom|autom.} d'ordre deux de \(X\), ou encore de} deux ``axes'' (ou diviseurs étales de degré \(2\)) \(\Delta_1,\Delta_2\) sur \(X\), la condition de commutation étant que \(\Delta_2\) est stable par \(\sigma_1\), ou \deleted{\unsure{en}} encore \(\Delta_1\) stable par \(\sigma_2\).
% slanted
\margin{on est en car. \(\ne2\)\\
\(\Delta_1=X^{\sigma_1}\)\\
\(\unsure{\Delta_2}=X^{\sigma_2}\)}
\deleted{Pour \(\sigma_1\) donné i.e. \(\Delta_1\) donné, \unsure{i.e. de \(\Delta_2\)} la donnée de \(\sigma_2\) \unsure{revient}} \inserted{La condition que (1) soit un mono signifie} \deleted{\unsure{que}} \unsure{aussi} que l'on ; \deleted{donc \unsure{à la donnée|que la donnée} d'une section de \(X/\{1,\sigma_1\}\unsure{\to|-}\)}
\[
\Delta_1\cap\Delta_2=\varnothing\qquad\text{i.e.}\quad\Delta_2\subset X-\Delta_1 \tag*{(2)}
\]
et pour \(\sigma_1\) \inserted{i.e. \(\Delta_1\)} donné, la donnée de \(\deleted{\unsure{e_2}}\ \sigma_2\) i.e. de \(\Delta_2\) équivaut à celle d'une section de \deleted{\unsure{\(X-\Delta_1\) sur|\(\Delta_1\) sur}} \((X-\Delta_1)/\{1,\sigma_1\}\) sur \(S\) (\(\Delta_2\) étant alors l'image inverse de cette section dans \(X-\Delta_1\)). Considérons le troisième larron \(e\)
\[
e_3=e_1+e_2 \tag*{(3)}
\]
il lui correspond un axe \(\Delta_3\), et \(\Delta_1,\Delta_2,\Delta_3\) sont mutuellement disjoints.
\endgroup
